%!TEX root = ../template.tex
%%%%%%%%%%%%%%%%%%%%%%%%%%%%%%%%%%%%%%%%%%%%%%%%%%%%%%%%%%%%%%%%%%%%
%% abstract-pt.tex
%% NOVA thesis document file
%%
%% Abstract in Portuguese
%%%%%%%%%%%%%%%%%%%%%%%%%%%%%%%%%%%%%%%%%%%%%%%%%%%%%%%%%%%%%%%%%%%%

\ExplSyntaxOn
\cs_if_exist:NT \TitleOnHeader
  { \TitleOnHeader }
\ExplSyntaxOff

\typeout{NT FILE abstract-pt.tex}%

Os Conflict-Free Replicated Data Types (CRDTs) representam uma abstração fundamental para sistemas distribuídos modernos, permitindo a atualização concorrente de dados em múltiplas réplicas e assegurando uma Consistência Eventual Forte sem coordenação centralizada. Apesar destas garantias, a construção correta de CRDTs continua a ser um desafio propenso a erros devido a poder comprometer a integridade dos dados. Embora as ferramentas de verificação automática baseadas em solvers SMT, permitam a criação rápida de protótipos e a validação de propriedades de convergência, estas enfrentam limitações significativas ao lidar com definições recursivas, propriedades indutivas e invariantes de estado global complexas, resultando frequentemente na impossibilidade de concluir a verificação.

Esta dissertação responde a estes desafios propondo uma framework integrada para o desenvolvimento modular e formal de CRDTs, que combina diversas ferramentas de backend baseadas em dedução. A framework tira partido da automação fornecida pelo VeriFx para verificar propriedades algébricas de convergência, enquanto recorre ao poder expressivo da plataforma Why3 para assegurar correções através de contratos e invariantes explícitos. É apresentada uma tradução formal das especificações VeriFx para WhyML, permitindo o uso de verificação dedutiva para impor uma semântica correta por construção onde a automação isolada é insuficiente. O trabalho preliminar, que incluiu a tradução e verificação de CRDTs como o Positive-Negative Counter tanto no VeriFx como no Why3, indica que as especificações podem ser partilhadas entre ferramentas sem perda de correção. A tradução revela as invariantes e propriedades estruturais que necessitam ser explicitas num ambiente dedutivo, motivando a integração de múltiplas técnicas de verificação.

% E agora vamos fazer um teste com uma quebra de linha no hífen a ver se a \LaTeX\ duplica o hífen na linha seguinte se usarmos \verb+"-+… em vez de \verb+-+.
%
% zzzz zzz zzzz zzz zzzz zzz zzzz zzz zzzz zzz zzzz zzz zzzz zzz zzzz zzz zzzz comentar"-lhe zzz zzzz zzz zzzz
%
% Sim!  Funciona! :)

% Palavras-chave do resumo em Português
% \begin{keywords}
% Palavra-chave 1, Palavra-chave 2, Palavra-chave 3, Palavra-chave 4
% \end{keywords}
\keywords{
  CRDTs \and
  Verificação Formal \and
  Why3 \and
  VeriFx
}
% to add an extra black line
